\documentclass[11pt]{article}
\usepackage[margin=1in]{geometry}
\usepackage{amsmath,amssymb,amsthm}
\usepackage[hidelinks]{hyperref}

\newtheorem{theorem}{Theorem}
\newtheorem{lemma}[theorem]{Lemma}
\newtheorem{proposition}[theorem]{Proposition}

\title{Extremal Heap-Type Permutation Generation for Every Odd $n$}
\author{Salem Levinstein}
\date{November 2026}

\begin{document}
\maketitle

\begin{abstract}
We answer in the affirmative an open question of Lu\'is M. S. Russo
(\texttt{arXiv:2611.22780v1}, Section~10, item~2) by proving that for every
odd integer $n \ge 3$, there exists a Heap-type single-swap enumeration of all
$n!$ permutations starting at $(1,2,\dots,n)$ and terminating at the
lexicographically largest odd permutation $a^*_n$. Only level $n-1$ departs
from Russo's cyclic schedule (using an explicit zigzag parking order), and both
the level-$(n-1)$ and top-level parking orders are given in closed form. The
entire construction and its end-to-end equivalence to the recursive swap
algorithm are formally verified in Lean~4.
\end{abstract}

\section{Introduction}

A \emph{Heap-type} permutation generator of order $n$ is specified by swap
schedules $s_k : \{1,\dots,k-1\} \to \{1,\dots,k-1\}$ for levels $2 \le k \le n$:
the procedure $\mathrm{Heap}(k)$ invokes $\mathrm{Heap}(k-1)$, then for
$m = 1,\dots,k-1$ swaps positions $s_k(m)$ and $k$ and invokes $\mathrm{Heap}(k-1)$
again. Starting from the identity array $(1,2,\dots,n)$, any valid Heap-type
generator visits all $n!$ permutations once, advancing by a single transposition
at each step. Because $n! - 1$ is odd for $n \ge 2$, the final arrangement must
be an odd permutation, so it is lexicographically bounded above by the largest
odd permutation $a^*_n$:
\[
a^*_n = \begin{cases}
(n,\, n-1,\, \dots,\, 2,\, 1) & \text{if } n \equiv 2,3 \pmod 4, \\
(n,\, n-1,\, \dots,\, 3,\, 1,\, 2) & \text{if } n \equiv 0,1 \pmod 4.
\end{cases}
\]

In~\cite{Russo2026}, Russo constructed a Heap-type generator reaching $a^*_n$ for
every $n \equiv 0 \pmod 4$ and showed that keeping levels $2,\dots,n-1$ cyclic
fails for odd $7 \le n \le 31$, leaving open (Section~10, item~2; Section~7
remark):
\begin{quote}
\emph{``We found no closed formula for the parking order at the top level and no
variant for odd $n \ge 7$ with the cyclic variant below it; whether other lower
levels make the largest attainable arrangement reachable for odd $n$ is open.''}
\end{quote}
Here we settle this question completely for every odd $n \ge 3$.

\begin{theorem}\label{thm:main}
For every odd integer $n \ge 3$, there exists a valid Heap-type swap schedule on
$\{1,\dots,n\}$ that starts at $(1,2,\dots,n)$, visits every permutation in
$S_n$ exactly once with one transposition per step, and terminates at $a^*_n$.
Both the level-$(n-1)$ and level-$n$ parking orders admit explicit closed forms.
\end{theorem}

Combined with Russo's construction for $n \equiv 0 \pmod 4$ and independent work
for $n \equiv 2 \pmod 4$, Theorem~\ref{thm:main} implies that the lexicographic
upper bound $a^*_n$ is attained by a Heap-type enumeration for every $n \ge 2$.

\section{Level Lifting via Parking Orders}

We adopt Russo's state convention: permutations compose right to left, an array
$a$ corresponds to the inverse state $g = a^{-1}$, and swapping positions $i$
and $j$ of $a$ multiplies $g$ on the left by $(i\;j)$.

\begin{lemma}[Level lifting]\label{lem:lift}
Suppose levels $1,\dots,k-1$ form a valid Heap-type enumeration of $S_{k-1}$
with net state $C \in S_{k-1}$ (fixing all positions $\ge k$). For any ordering
$(e_0, e_1, \dots, e_{k-1})$ of $\{1,\dots,k\}$ with $e_0 = k$, there exists a
level-$k$ swap schedule $s_k : \{1,\dots,k-1\} \to \{1,\dots,k-1\}$ such that
$\mathrm{Heap}(k)$ is valid and has net state
\[
\pi_k = C^k (k\; e_1\; e_2\; \cdots\; e_{k-1}).
\]
\end{lemma}

\begin{proof}
Define block-initial states inductively by $g_0 = \mathrm{id}$ and, for
$0 \le m < k-1$, set $h_m = C g_m$, $s_k(m+1) = h_m(e_{m+1})$, and
$g_{m+1} = (s_k(m+1)\;k)\,h_m$. By induction on $m$, $g_m(e_m) = k$, whence
$h_m(e_m) = C(k) = k$. Since $e_{m+1} \neq e_m$ and $h_m \in S_k$, we have
$s_k(m+1) \in \{1,\dots,k-1\}$, so every swap is in range. Conjugating by $h_m$
gives
\[
g_{m+1} = (h_m(e_{m+1})\; h_m(e_m))\, h_m = h_m (e_{m+1}\;e_m) = C g_m (e_{m+1}\;e_m),
\]
and therefore $g_m = C^m (e_1\;e_0)(e_2\;e_1)\cdots(e_m\;e_{m-1})$. Every state
$p g_m$ visited during block $m$ satisfies $(p g_m)^{-1}(k) = g_m^{-1}(k) = e_m$.
Because $e_0,\dots,e_{k-1}$ are pairwise distinct, the $k$ blocks are disjoint
cosets of $S_{k-1}$, so $\mathrm{Heap}(k)$ visits all $k!$ permutations once.
Finally, its net state is
\[
\pi_k = C g_{k-1} = C^k (e_1\;e_0)(e_2\;e_1)\cdots(e_{k-1}\;e_{k-2})
= C^k (k\; e_1\; e_2\; \cdots\; e_{k-1}). \qedhere
\]
\end{proof}

\begin{lemma}[Cyclic lower levels]\label{lem:cyclic}
For every $k \ge 2$, there is a valid Heap-type schedule through level $k$ whose
net state is the $k$-cycle $(k\;1\;2\;\cdots\;k-1)$ when $k$ is even and the
transposition $(1\;k)$ when $k$ is odd.
\end{lemma}

\begin{proof}
Induct on $k$. For $k=2$, $C = \mathrm{id}$ and $e_1 = 1$ give net $(2\;1)$.
For even $k \ge 4$, the previous odd net $C = (1\;k-1)$ satisfies $C^k = \mathrm{id}$,
so Lemma~\ref{lem:lift} with $(e_1,\dots,e_{k-1}) = (1,2,\dots,k-1)$ gives
$(k\;1\;2\;\cdots\;k-1)$. For odd $k \ge 3$, the previous even net
$C = (k-1\;1\;2\;\cdots\;k-2)$ has order $k-1$, so $C^k = C$; applying
Lemma~\ref{lem:lift} with descending order $(e_1,\dots,e_{k-1}) = (k-1,k-2,\dots,1)$
gives $C (k\;k-1\;\cdots\;1) = (1\;k)$.
\end{proof}

\section{Proof of Theorem~\ref{thm:main}}

For $n = 3$, Lemma~\ref{lem:cyclic} yields a valid Heap-type schedule with net
$(1\;3)$, whose inverse array is $(3,2,1) = a^*_3$.

Now assume $n \ge 5$ is odd, and set $h = (n - 3)/2 \ge 1$.
\begin{enumerate}
\item \textbf{Levels $2,\dots,n-2$.} Apply Lemma~\ref{lem:cyclic} through the odd
level $n - 2$, obtaining net $C = (1\;n-2)$.

\item \textbf{Level $n - 1$ (zigzag parking order).} Because $C = (1\;n-2)$ has
order $2$ and $n - 1$ is even, $C^{n-1} = \mathrm{id}$. Apply
Lemma~\ref{lem:lift} at level $n - 1$ with the zigzag parking order
\[
E' = (1,\; n-2,\; 2,\; n-3,\; \dots,\; h,\; h+2,\; h+1).
\]
Its net is the $(n - 1)$-cycle $Z' = (n - 1\; e'_1\;\cdots\; e'_{n-2})$, given
explicitly by
\[
Z'(x) = \begin{cases}
n - 1 - x & 1 \le x \le h, \\
n - 1 & x = h + 1, \\
n - x & h + 2 \le x \le n - 1, \\
n & x = n,
\end{cases}
\qquad
(Z')^{-1}(y) = \begin{cases}
n - y & 1 \le y \le h + 1, \\
n - 1 - y & h + 2 \le y \le n - 2, \\
h + 1 & y = n - 1, \\
n & y = n.
\end{cases}
\]
Since $Z'$ is an $(n - 1)$-cycle, $(Z')^{n-1} = \mathrm{id}$ and hence $(Z')^n = Z'$.

\item \textbf{Level $n$ (three-phase top order).} The target state is
$T = (a^*_n)^{-1}$, which satisfies $T(v) = n + 1 - v$ for $3 \le v \le n$,
$(T(1), T(2)) = (n, n - 1)$ if $n \equiv 3 \pmod 4$, and
$(T(1), T(2)) = (n - 1, n)$ if $n \equiv 1 \pmod 4$. Set $Z = (Z')^{-1} T$.
Evaluating $(Z')^{-1}(T(v))$ directly yields:
\begin{itemize}
\item $Z(v) = v - 2$ for $3 \le v \le h + 2$ (as $T(v) \in [h + 2, n - 2]$);
\item $Z(v) = v - 1$ for $h + 3 \le v \le n$ (as $T(v) \in [1, h + 1]$);
\item $(Z(1), Z(2)) = (n, h + 1)$ if $n \equiv 3 \pmod 4$, and
$(Z(1), Z(2)) = (h + 1, n)$ if $n \equiv 1 \pmod 4$.
\end{itemize}
\end{enumerate}

We verify that $Z$ is a single $n$-cycle by following the orbit of $n$:
\begin{itemize}
\item \emph{Phase~1:} $n \mapsto n - 1 \mapsto \dots \mapsto h + 2$ (decreasing by $1$).
\item \emph{Phase~2:} from $h + 2 = (n + 1)/2$, decreasing by $2$ traverses every
integer in $[1, h + 2]$ of the same parity as $h + 2$, reaching $r \in \{1, 2\}$
with $r \equiv h + 2 \pmod 2$. If $n \equiv 3 \pmod 4$ then $h + 2$ is even,
$r = 2$, and $Z(2) = h + 1$; if $n \equiv 1 \pmod 4$ then $h + 2$ is odd,
$r = 1$, and $Z(1) = h + 1$.
\item \emph{Phase~3:} from $h + 1 = (n - 1)/2$ (the opposite parity), decreasing
by $2$ traverses every integer in $[1, h + 1]$ of the parity of $h + 1$, reaching
$3 - r \in \{1, 2\}$, where $Z(3 - r) = n$.
\end{itemize}
The three phases cover $[h + 2, n]$ and both parity classes of $[1, h + 1]$
disjointly, so $Z$ is an $n$-cycle with explicit top-level parking order
\[
E = (n - 1,\, n - 2,\, \dots,\, h + 2,\; h,\, h - 2,\, \dots,\, r,\; h + 1,\, h - 1,\, \dots,\, 3 - r).
\]
Applying Lemma~\ref{lem:lift} at level $n$ with order $E$ produces a valid
Heap-type schedule with net $\pi_n = (Z')^n Z = Z' (Z')^{-1} T = T$, hence final
array $T^{-1} = a^*_n$. \qed

\section{Formal and Computational Verification}

The entire proof---including the recursive definition of Heap-type enumerations,
Lemmas~\ref{lem:lift} and~\ref{lem:cyclic}, the sign and lexicographic
maximality of $a^*_n$, and Theorem~\ref{thm:main}---is formally verified in
Lean~4 (\texttt{Conjectures/RussoHeapOddExtremal.lean}, theorem
\texttt{russo\_open\_question\_item2}) with zero \texttt{sorry} and standard
axioms only. An independent Python script (\texttt{verify.py}) directly
simulates the recursion for $n \in \{3, 5, 7, 9\}$ and verifies all closed-form
identities for every odd $5 \le n \le 1201$.

\begin{thebibliography}{9}
\bibitem{Russo2026}
L.~M.~S. Russo,
\emph{Heap-type permutation generation with a prescribed final arrangement},
arXiv:2611.22780v1, November 2026.
\end{thebibliography}

\end{document}
